\section{Background}\label{sec:background}

\begin{comment}
	{\bf Write-ahead logging in databases.} Write-ahead logging (WAL) is a fundamental technique that databases rely on for consistency and durability~\cite{mohan1992aries,gray1993transaction}.
	Databases persistently write key-value pairs or operations to log files for recovery in case of power outage or software panic.
	As log files are rarely read except during recovery, they typically operate in direct I/O mode, bypassing the OS page cache.
	Database systems employ two primary strategies for log file management: preallocation and dynamic growth.
	Rising-star databases such as OceanBase initialize log files by preallocating tens of gigabytes by default~\cite{DB:OceanBase:VLDB-2022}.
	Similarly, mature database engines like SQLite can be tuned to preallocate log files~\cite{lee2015waldio}.
	However, many production systems also use dynamically growing log files, where blocks are allocated on-demand as the log expands, presenting unique performance challenges.
\end{comment}

%{\bf The I/O path anatomy: where software tax accumulates.}
\textbf{File I/O Path.}
A file operation that an application issues %, such as a file write or fdatasync,
%When an application issues a file operation, the I/O request 
traverses a complex software stack composed of multiple layers 
before reaching the storage device.
The top  layer is  {\em syscall interface}, at which syscalls such as file write and fdatasync 
transition to the kernel space with parameter validation.
The {\em VFS} then performs inode lookup, interacts with the OS page cache, 
manages file position tracking, and so on~\cite{love2010linux}.
%The file operation is next passed to 
VFS passes the file operation to
a {\em specific file system}, such as the prevalent Ext4 studied in this paper.
Ext4
maps   in-file offsets to device block numbers and   composes block I/O (bio) requests.
For a write issued to unallocated space, Ext4 allocates   space on demand
and updates the per-file mappings. These metadata changes will trigger
{\em journaling} for crash consistency, by committing modified metadata
%are committed 
to a dedicated {\em journal}
before in-place update~\cite{mathur2007ext4,fsync:FastCommit:ATC-2024,fsync:iJournaling:ATC-2017}.
The {\em bio layer} and {\em I/O scheduler}  manage, merge, and schedule bio requests where necessary%
%bio requests, merge requests , and schedule 
%requests
~\cite{love2010linux}.
Finally, the {\em device driver} translates bio requests to device-specific commands. It collects completion or fail signals from the device and
notifies upper-level  layers.

 

\textbf{Software Tax.}	
Software layers impose heavy tax on the critical path of serving a file write.
%When traversing these  layers, 	software overhead will occur on the critical path.
In particular, locating a target block for a write involves a lookup
in the file's extent tree, a structure by which Ext4 stores the  mappings per file. 
%The block allocation triggers %, which
The journaling caused by allocating blocks and updating the per-file mappings,
 entails severe write amplification with multiple writes performed to  the journal and file \cite{lee2015waldio,fsync:iJournaling:ATC-2017,fsync:FastCommit:ATC-2024}.
%checking the per-file mappings involves a lookup in an extent tree
 
%In the device, OS kernel, and applications, researchers have incorporated a few tactics to accelerate file writes.
Researchers have considered accelerating %incorporated a few tactics to accelerate
 file writes in the device, OS, and applications.
Firstly, many enterprise-level SSDs are equipped with the supercapacitor-based  PLP,
% technique
% that enables a temporary, 
which flushes 
%, thereby enabling
%PLP allows 
%the flush of
%PLP allows the SSD 
data %or metadata 
staying in the SSD's internal cache to be durable in case of power outage \cite{PLP:samsung}.
%in case of power outage, thereby dramatically
Calling fsync to write data into one such SSD can return promptly 
once the data arrives at the cache.  PLP hence largely shortens the response time of fsync%
% thereby largely reducing the response time of fsync%
%This largely reduces the response time of fsync/fdatasync
~\cite{fsync:barrierFS:FAST-2018}.
Secondly, in the OS kernel,
% researchers have explored at multiple aspects %dimensions 
%in the OS. For example, 
instead of using
the stricter {\em mandatory access control} (MAC) to check permission   % to check the access permission 
at every file operation,
the OS mainly applies the  {\em discretionary access control} (DAC) only at opening a file. 
%For buffered writes, Ext4 provides %file systems like Ext4 use 
% the delayed block allocation option. % for device blocks.
%VFS and file systems maintain dentry (directory entry) cache and inode cache for acceleration.
As for persistent stores,
the OS provides  fatasync to substitute fsync if there is no metadata change to a file.
% In case of no metadata change,  
%Researchers  have optimized the journaling mechanism   also
The journaling mechanism has been optimized also \cite{fsync:RFLUSH:FAST-2018,fsync:rethink:OSDI-2006,FS:Z-Journal:ATC-2021,fsync:OptFS:SOSP-2013,fsync:TxFS:ATC-2018,fsync:FastCommit:ATC-2024}. 
Thirdly,  developers have practical methods in their applications.
To minimize the cost of on-demand space allocations and 
metadata journaling,   many databases
preallocate all or partial files in advance \cite{DB:OceanBase:VLDB-2022,MySQL:fallocate,SQL-server:fallocate}.  %, such as OceanBase.
Preallocation
 also facilitates the use of DAC for permission checks   and
fdatasync for durable stores.
%the caches of inodes and directory entries (dentries)

%complete critical operations 
%in case of power outage.




\begin{comment}
	
	
	\textit{1. System call interface:} The initial entry point where user-space transitions to kernel-space through syscalls like write(), read(), or fsync(). This involves context switching overhead and parameter validation.
	
	\textit{2. Virtual File System (VFS):} The abstraction layer that provides a uniform interface for different file systems. VFS performs inode lookups, maintains file descriptors, and manages file position tracking~\cite{love2010linux}.
	
	\textit{3. File system layer (e.g., Ext4):} This layer performs the critical task of translating logical file offsets to physical disk blocks.
	For Ext4, this involves navigating extent trees—B+ tree structures that map file regions to contiguous disk blocks~\cite{mathur2007ext4}.
	Dynamic block allocation adds substantial overhead: allocating new blocks requires searching free space bitmaps, updating block groups, and potentially triggering expensive journal commits~\cite{prabhakaran2005analysis}.
	
	\textit{4. Journaling subsystem:} Modern file systems like Ext4 use journaling (JBD2) to ensure consistency~\cite{tweedie1998journaling}.
	Each metadata modification—including block allocations for growing files—must be logged to the journal before the actual update.
	This creates a write amplification effect: a single logical write can trigger multiple physical writes to both the journal and the file data.
	
	\textit{5. Block layer and I/O scheduler:} The generic block layer manages I/O requests, performs request merging, and applies scheduling policies~\cite{love2010linux}.
	While beneficial for throughput, these optimizations add latency for latency-sensitive logging operations.
	
	\textit{6. Device driver:} Finally, the device driver translates block I/O requests into device-specific commands (e.g., NVMe commands) and manages completion interrupts.
	
	
	{\bf Permission checks: DAC vs. MAC.}
	Linux implements two distinct access control models that affect I/O performance differently~\cite{smalley2001implementing,wright2002linux}:
	
	\textit{Discretionary Access Control (DAC):} The traditional Unix permission model checks file access rights (read/write/execute) based on user/group/other permissions.
	Critically, DAC performs these checks only once when a file is opened, storing the validated permissions in the file descriptor.
	Subsequent I/O operations through that descriptor bypass permission checks entirely.
	
	\textit{Mandatory Access Control (MAC):} Security frameworks like SELinux and AppArmor implement MAC policies that can enforce fine-grained access controls~\cite{loscocco2001integrating}.
	Unlike DAC, MAC systems may perform access checks on every I/O operation, not just at file open time.
	This recurring validation ensures stronger security but introduces measurable overhead—each write() or read() syscall triggers a policy evaluation that can involve complex rule matching.
\end{comment}




% huyp: have been mentioned in the intro
\begin{comment}
	Emerging memory technologies, like Intel 3D XPoint \cite{intel3dxpoint} and low-latency NAND flash \cite{samsungznand2024}, have paved the way for NVMe SSDs that offer $\mu$s-level latency and millions of I/O operations per second (IOPS).
	With increasingly faster devices, the tax imposed by software layers turns to rise in the overall time cost per I/O request. 
	
	{\bf SOTA solutions.}
	Researchers have studied the reduction or even avoidance of software tax.
	State-of-the-art (SOTA) solutions include NVMeDirect \cite{SSD:NVMeDirect:HotStorage-2016}, SPDK \cite{SSD:SPDK:CloudCom-2017}, Moneta-D \cite{SSD:Moneta-D:ASPLOS-2012}, BypassD \cite{SSD:BypassD:ASPLOS-2024}, etc.
	Using the interfaces provided by NVMeDirect or SPDK, developers can fully bypass the OS kernel. Though, they must reprogram applications and manage data on their own.
	Moneta-D and BypassD incorporate customized hardware for a part of management purposes. In particular, Moneta-D mainly makes a specialized storage device perform permission checks.
	BypassD extends the IOMMU to map in-file offsets to disk blocks and check block-level access permissions.
	Both of them also need to change the file system in the OS kernel.
\end{comment}



\textbf{eBPF.}
The extended Berkeley Packet Filters (eBPF) \cite{ebpf} was initially developed to filter packets in networks. It now enables advanced capabilities like performance monitoring, security enforcement, and I/O optimization by allowing custom code to run efficiently in the OS kernel.
A notable feature of eBPF is its ability to transparently alter the execution path of a kernel function and run alternate codes that have been verified \cite{ebpf,280870,DBLP:conf/osdi/ZhengY0HL0Q25}.
This makes a seamless transition to a new execution path for a particular syscall, without requiring any change to applications.
eBPF also supports a memory buffer shared at the boundary of kernel and user spaces.
We will exploit these features to exempt the heavy runtime tax charged by the software layers on
file writes.
% on logging I/Os.


\textbf{ioctl.} The ioctl ({\em I/O control}) enables  user-space applications
to directly interact with a device \cite{ioctlWik84:online}. If the numbers of device blocks
in which a file's data is stored are known, 
a program can leverage the ioctl interface to directly write data into these blocks
and   flush the device's internal cache where necessary\cite{6702624}.
Though, the ioctl interfaces must be used with particular care. For example, 
an incorrect block number may mislead an overwrite to 
some file's inode or the
file system's metadata (e.g., superblock), corrupting 
a file or even the entire system.
The aforementioned preallocation helps to obtain the numbers of   
blocks  belonging to a file.
%Given file data stored in some block numbers,

%TODO: 